% Generated draft: requires mathematical and source review.
\documentclass[11pt]{article}
\usepackage[margin=25mm]{geometry}
\usepackage{fontspec}
\setmainfont{Latin Modern Roman}
\usepackage{amsmath,amssymb,mathtools,unicode-math}
\setmathfont{Latin Modern Math}
\usepackage{enumitem,needspace,xurl,hyperref}
\setlength{\parindent}{0pt}
\setlength{\parskip}{5pt}
\setlength{\emergencystretch}{3em}
\newcommand{\N}{\mathbb{N}}
\newcommand{\Z}{\mathbb{Z}}
\newcommand{\Q}{\mathbb{Q}}
\newcommand{\R}{\mathbb{R}}
\newcommand{\C}{\mathbb{C}}
\newcommand{\F}{\mathbb{F}}
\newcommand{\A}{\mathbb{A}}
\newcommand{\PP}{\mathbb{P}}
\newcommand{\E}{\mathbb{E}}
\newcommand{\eps}{\varepsilon}
\newcommand{\dd}{\,\mathrm{d}}
\begin{document}
\section*{UnsolvedMath ID: 800491}
\textbf{Source problem code:} catalogue entry
\subsection{Definitions and mathematical statement}
Let \(\mathbf{C}P^n\) denote complex projective space. Relative to a symplectic form \(\omega\), a Hamiltonian diffeomorphism \(\varphi\) is the time-one map of a smooth time-dependent Hamiltonian flow. The supplied record does not specify \(\omega\) or explicitly quantify the dimension \(n\).

A periodic point of \(\varphi\) is a point \(x\) such that \(\varphi^k(x)=x\) for some positive integer \(k\). Its least period is
\[
\tau(x)=\min\{k\in\mathbb{Z}_{\geq 1}:\varphi^k(x)=x\}.
\]
Its periodic orbit is the finite set \(\mathcal O_\varphi(x)=\{x,\varphi(x),\ldots,\varphi^{\tau(x)-1}(x)\}\). Thus distinct points can belong to the same periodic orbit.

The supplied question asks whether every Hamiltonian diffeomorphism of \(\mathbf{C}P^n\) with infinitely many periodic orbits satisfies
\[
\liminf_{T\to+\infty}\frac{\log P(T)}{\log T}\geq 2.
\]
Here \(P(T)\) is not defined in the supplied record. Two possible counting functions, for real \(T>1\), are
\[
P_{\mathrm{orb}}(T)=\#\{\mathcal O_\varphi(x):x\text{ is periodic and }\tau(x)\leq T\},
\qquad
P_{\mathrm{pt}}(T)=\#\{x:x\text{ is periodic and }\tau(x)\leq T\}.
\]
The first counts each orbit once; the second counts its points separately. These are alternative interpretations, not definitions supplied by the source. The question uses periodic orbits, whereas the discussion of known bounds uses periodic points. The intended choice of \(P(T)\), the intended period convention, and the treatment of infinite counts therefore require clarification. No assumption ensuring finite counts at each bounded period is supplied. In particular, the displayed logarithmic expression is unresolved when a count is infinite unless an extended-real convention is adopted.

The source motivates the exponent \(2\) through low-dimensional integrable systems. It describes examples with exactly quadratic growth when one counts invariant tori foliated by periodic orbits together with isolated closed geodesics on a closed Riemannian \(2\)-sphere of revolution. This is a modified counting problem, and the source does not define its counting function or the precise meaning of exact quadratic growth. It conjectures that perturbations could produce systems with exact quadratic growth of periodic orbits, while emphasizing that constructing such perturbations may be very difficult. It also reports an analogous expectation for integrable annulus maps, attributed to Gerhard Knieper. These observations do not supply additional hypotheses or an established sharpness result for the displayed question.

For higher-dimensional motivation, the source recalls the Hofer--Zehnder question [HZ, p.\ 263]: whether a Hamiltonian map on a compact symplectic manifold \((M,\omega)\) with more fixed points than required by the Arnold conjecture must have infinitely many periodic orbits. The precise fixed-point threshold and its hypotheses are not specified in the record. The source attributes to Shelukhin [Shelukhin] a proof for closed monotone symplectic manifolds whose even-dimensional quantum homology ring is semisimple, a class stated to include complex projective spaces. It reports that this proof yields
\[
\liminf_{T\to+\infty}P(T)\frac{\log T}{T}>0,
\]
with \(P(T)\) described there as counting periodic points, and compares this growth bound with [Hingston1993]. These are supplied literature claims; the undefined counting convention and omitted theorem hypotheses prevent treating them here as fully specified theorem statements.

The source further attributes to Ginzburg and G\"urel a broader interpretation in which a fixed point not strictly required by Floer theory forces infinitely many periodic points. Its examples are a hyperbolic fixed point [Ginzburg-Gurel2014] and a non-contractible fixed point [Gurel2013, Ginzburg-Gurel2016]. The ambient hypotheses and the meaning of non-contractibility relative to a Hamiltonian isotopy are not supplied. These are broader contextual examples, not additional assumptions on the projective-space question. The source also attributes the preceding positive lower bound to G\"urel [Gurel2013] and tentatively asks whether quadratic growth holds in such settings as well; it gives no precise common formulation of that extension.

The registry labels this problem open. That label and the literature descriptions above are supplied source data, not independently verified mathematical assessments.
\subsection{Short English statement}
The source asks whether every Hamiltonian diffeomorphism of \(\mathbf{C}P^n\) with infinitely many periodic orbits has lower logarithmic growth exponent at least \(2\). A precise formulation still requires defining \(P(T)\), distinguishing periodic points from periodic orbits, and specifying how infinite counts are treated. The suggested sharpness examples and broader extensions remain qualified source remarks.
\subsection{Sources}
\begin{itemize}
\item Local UnsolvedMath registry, permanent ID \(800491\), exact source code \texttt{catalogue entry}, AMR Open Problem Lists. The supplied record has status \texttt{open}, proposed year \(2026\), and collection-level attribution to Albach, Barraud et al.
\item Albach, Barraud et al., \emph{Open Problems in Billiards and Quantitative Symplectic Geometry} (2026), \url{https://arxiv.org/abs/2602.12896}. The supplied title and provenance identify Bernhard Albach, Umberto Hryniewicz, and Alexandru Oancea, Problem 4; the extracted statement instead begins ``Problem 3.'' The same three people are listed as contributors. Supplied access date: 2026-10-16; extraction method: \texttt{arxiv\_latex}. This date is supplied provenance, not a claim of independent access or verification.
\item [HZ, p.\ 263]: supplied attribution to Hofer and Zehnder for the fixed-point question. No complete bibliographic entry is included.
\item [Shelukhin]: supplied attribution for the stated Hofer--Zehnder result and periodic-point growth estimate. No complete bibliographic entry is included.
\item [Hingston1993]: supplied comparison for the periodic-point growth estimate. No complete bibliographic entry is included.
\item [Ginzburg-Gurel2014]: supplied attribution to Ginzburg and G\"urel concerning a hyperbolic fixed point. No complete bibliographic entry is included.
\item [Gurel2013] and [Ginzburg-Gurel2016]: supplied references concerning the non-contractible case; the growth estimate is also attributed to [Gurel2013]. No complete bibliographic entries are included.
\item Gerhard Knieper: supplied attribution for the expectation concerning integrable annulus maps; no separate publication or URL is supplied.
\end{itemize}
\end{document}
